\section{Table-Form Kronecker Encoding}
\label{sec:table-encoding}

The protocols in this paper commit directly to vectors indexed by the Boolean hypercube.  The corresponding univariate polynomial is obtained by using the integer index of a Boolean point as the exponent.  This section defines that encoding, relates it precisely to the canonical and Lagrange representations of multilinear polynomials from \cref{sec:preliminaries-multilinear}, and records the commitment properties used throughout the paper.

\subsection{The table polynomial}
\label{sec:table-polynomial}

Let $t:\B^n\to\F$ be a table.  As in \cref{sec:preliminaries-bits}, we identify a Boolean point
\[
  w=(w_1,\ldots,w_n)\in\B^n
\]
with the integer
\[
  w=\sum_{i=1}^{n}w_i2^{i-1}\in[0,N-1].
\]
Thus we write $t(w)$ interchangeably for the value at the bit vector and at its integer index.

\begin{definition}[Table polynomial]
\label{def:table-polynomial}
For a table $t:\B^n\to\F$, its \emph{table polynomial} is
\begin{equation}
  V_t(X)
  \defeq
  \sum_{w=0}^{N-1}t(w)X^w
  \in\F[X]_{<N}.
  \label{eq:table-polynomial}
\end{equation}
\end{definition}

The construction in \cref{def:table-polynomial} is the Kronecker encoding of the length-$N$ vector $(t(0),\ldots,t(N-1))$: the entries of the vector are used verbatim as the coefficients of a univariate polynomial.  No interpolation on the Boolean cube is involved in the definition of $V_t$.

\begin{proposition}[Table encoding is a linear isomorphism]
\label{prop:table-isomorphism}
The map
\[
  \mathcal K_T:\F^{\B^n}\longrightarrow\F[X]_{<N},
  \qquad
  t\longmapsto V_t,
\]
is an $\F$-linear isomorphism.  Its inverse maps
\[
  P(X)=\sum_{j=0}^{N-1}p_jX^j
\]
to the table $t_P$ defined by $t_P(j)=p_j$.
\end{proposition}

\begin{proof}
Linearity follows coefficientwise from \eqref{eq:table-polynomial}.  Every polynomial in $\F[X]_{<N}$ has a unique coefficient vector $(p_0,\ldots,p_{N-1})$, and therefore a unique table $t_P$ with $t_P(j)=p_j$.  For this table,
\[
  V_{t_P}(X)=\sum_{j=0}^{N-1}p_jX^j=P(X).
\]
Conversely, the coefficient of $X^j$ in $V_t$ is exactly $t(j)$, so $t_{V_t}=t$.  Thus the two maps are inverse linear maps.
\end{proof}

\begin{remark}[No distinction between a vector and its table polynomial]
\label{rem:vector-table-polynomial}
Whenever a vector $a=(a_0,\ldots,a_{N-1})\in\F^N$ is indexed by $[0,N-1]$, we may regard it as the table $a:\B^n\to\F$ under the fixed binary indexing and write
\[
  V_a(X)=\sum_{j=0}^{N-1}a_jX^j.
\]
This convention will be used for witness vectors, lookup columns, matrix-derived vectors, and intermediate pointwise-arithmetic wires.  The encoding depends only on the ordered vector, not on an a priori multivariate polynomial.
\end{remark}

\subsection{Relation with multilinear polynomial representations}
\label{sec:table-vs-multilinear}

The table polynomial $V_t$ should not be confused with the canonical-coordinate Kronecker polynomial of the multilinear extension $\widetilde t$.  Both are degree-$<N$ univariate polynomials, but their coefficients represent different coordinate systems.

Recall from \cref{prop:boolean-interpolation} that every table $t:\B^n\to\F$ has a unique multilinear extension
\[
  \widetilde t(z)
  =\sum_{w\in\B^n}t(w)\eq(w,z).
\]
The family
\begin{equation}
  e_w(x_1,\ldots,x_n)
  \defeq
  \eq(w,x)
  =\prod_{i=1}^{n}\bigl((1-w_i)(1-x_i)+w_ix_i\bigr)
  \label{eq:boolean-lagrange-basis}
\end{equation}
is the Boolean Lagrange basis of the multilinear space over $\F$.

\begin{proposition}[Table values are Lagrange coordinates]
\label{prop:table-lagrange-coordinates}
For every table $t:\B^n\to\F$,
\begin{equation}
  \widetilde t(x)
  =\sum_{w\in\B^n}t(w)e_w(x).
  \label{eq:mle-lagrange-expansion}
\end{equation}
Consequently, the coordinate vector of $\widetilde t$ in the Boolean Lagrange basis $(e_w)_{w\in\B^n}$ is exactly the table vector $(t(w))_{w\in\B^n}$, and $V_t$ is the univariate Kronecker encoding of those Lagrange coordinates.
\end{proposition}

\begin{proof}
Equation \eqref{eq:mle-lagrange-expansion} is the definition of the multilinear extension in \eqref{eq:mle-definition}.  The Boolean delta property \cref{lem:eq-delta} implies that the polynomials $e_w$ interpolate the standard coordinate vectors on $\B^n$, hence form the Lagrange basis.  Therefore the coefficient of $e_w$ in \eqref{eq:mle-lagrange-expansion} is $t(w)$.  By \cref{def:table-polynomial}, these same coordinates are the coefficients of $V_t$.
\end{proof}

Now let $f\in\M$ be written in the canonical basis of \cref{prop:canonical-basis} as
\[
  f=\sum_{j=0}^{N-1}\alpha_jm_j.
\]
Its canonical-coordinate Kronecker polynomial is
\begin{equation}
  U_f(X)
  \defeq
  \sum_{j=0}^{N-1}\alpha_jX^j.
  \label{eq:canonical-kronecker-polynomial}
\end{equation}
Its table polynomial, by contrast, is
\[
  V_f(X)
  \defeq
  \sum_{b=0}^{N-1}f(b)X^b.
\]
The two coefficient vectors are related by the Boolean zeta and M\"obius transforms from \cref{lem:mobius}.

\begin{proposition}[Canonical and table Kronecker polynomials]
\label{prop:canonical-table-kronecker}
Let $f\in\M$ have canonical coordinates $(\alpha_j)_{j=0}^{N-1}$.  Then
\begin{equation}
  [X^b]V_f
  =f(b)
  =\sum_{j\preceq b}\alpha_j,
  \label{eq:table-from-canonical}
\end{equation}
and
\begin{equation}
  [X^j]U_f
  =\alpha_j
  =\sum_{c\preceq j}(-1)^{\wt(j)-\wt(c)}[X^c]V_f.
  \label{eq:canonical-from-table}
\end{equation}
In particular, $U_f$ and $V_f$ generally differ even though both belong to $\F[X]_{<N}$.
\end{proposition}

\begin{proof}
The first equality in \eqref{eq:table-from-canonical} is the definition of $V_f$, and the second is \eqref{eq:zeta-transform}.  Equation \eqref{eq:canonical-from-table} is \eqref{eq:mobius-transform} after substituting $f(c)=[X^c]V_f$.  No further argument is required.
\end{proof}

\begin{remark}[Why table form is the default in this paper]
\label{rem:why-table-form}
For a protocol whose witness is already represented by values on $\B^n$, the coefficients of $V_t$ are available directly.  In particular, committing to $V_t$ does not require first applying the M\"obius transform \eqref{eq:mobius-transform}.  More importantly, later relations such as pointwise products, lookup columns, permutations, and sparse-matrix streams are naturally relations between table entries.  The table polynomial therefore gives one uniform univariate representation for all of them.
\end{remark}

\subsection{The Reed--Solomon table commitment}
\label{sec:table-commitment}

Fix the common parameters of \cref{sec:preliminaries-parameters}: a smooth multiplicative domain $L\le\Fq^\times$ of order $M\ge N$ and the base code
\[
  C_0=\RS_{\F}[L,N].
\]

\begin{definition}[Table-form Kronecker commitment]
\label{def:table-commitment}
For a table $t:\B^n\to\F$, define its oracle commitment by
\begin{equation}
  \Enc_T(t)
  \defeq
  \ev_L(V_t)
  =\bigl(V_t(\xi)\bigr)_{\xi\in L}
  \in C_0.
  \label{eq:table-commitment}
\end{equation}
In the interactive-oracle model, the word $\Enc_T(t)$ is the commitment oracle.  A later non-interactive compilation will replace the oracle by a hash/Merkle commitment to its entries.
\end{definition}

\begin{proposition}[Commitment is a linear bijection onto the base code]
\label{prop:table-commitment-bijection}
The map
\[
  \Enc_T:\F^{\B^n}\longrightarrow C_0
\]
is an $\F$-linear bijection.
\end{proposition}

\begin{proof}
By \cref{prop:table-isomorphism}, $t\mapsto V_t$ is an $\F$-linear bijection from $\F^{\B^n}$ onto $\F[X]_{<N}$.  By \cref{prop:rs-parameters}(1), evaluation on $L$ is an $\F$-linear bijection from $\F[X]_{<N}$ onto $C_0$.  Their composition is therefore an $\F$-linear bijection.
\end{proof}

The previous proposition gives a canonical decoder from an exact codeword to a table: if $c\in C_0$, let $P_c\in\F[X]_{<N}$ be its unique representing polynomial and define
\begin{equation}
  \operatorname{Tab}(c)(j)
  \defeq
  [X^j]P_c.
  \label{eq:codeword-table-decoder}
\end{equation}
Then $\Enc_T(\operatorname{Tab}(c))=c$.

For soundness we need the stronger statement that a word sufficiently close to the code determines at most one committed table.

\begin{proposition}[Unique table determined by a nearby word]
\label{prop:unique-table-near-word}
Let $0\le\delta\le\delta_*=(1-\rho)/2$, and let $w\in\F^L$.  There is at most one table $t:\B^n\to\F$ satisfying either of the following conditions:
\begin{align}
  \Delta\bigl(w,\Enc_T(t)\bigr)&\le\delta,
  \label{eq:table-near-hamming}\\
  \Delta_{\mathrm{fib}}\bigl(w,\Enc_T(t)\bigr)&\le\delta.
  \label{eq:table-near-fibre}
\end{align}
\end{proposition}

\begin{proof}
For \eqref{eq:table-near-hamming}, \cref{prop:rs-parameters}(3) implies that there is at most one codeword $c\in C_0$ within distance $\delta$ of $w$.  By \cref{prop:table-commitment-bijection}, each codeword in $C_0$ corresponds to exactly one table, so at most one such $t$ exists.

For \eqref{eq:table-near-fibre}, \cref{lem:hamming-fibre} gives
\[
  \Delta\bigl(w,\Enc_T(t)\bigr)
  \le
  \Delta_{\mathrm{fib}}\bigl(w,\Enc_T(t)\bigr)
  \le\delta.
\]
Thus any table satisfying the fibre-distance condition also satisfies the Hamming-distance condition, and uniqueness follows from the first part.
\end{proof}

\begin{proposition}[Base-field consistency]
\label{prop:table-basefield-consistency}
If $t:\B^n\to\Fq$, then
\[
  \Enc_T(t)\in\Fq^L.
\]
Conversely, let $w\in\Fq^L$ and suppose that for some table $t:\B^n\to\F$,
\begin{equation}
  \Delta\bigl(w,\Enc_T(t)\bigr)
  <\frac12\left(1-\frac{N-1}{M}\right).
  \label{eq:table-basefield-radius}
\end{equation}
Then $t$ takes values in $\Fq$.
\end{proposition}

\begin{proof}
If $t(w)\in\Fq$ for every $w$, then all coefficients of $V_t$ lie in $\Fq$.  Since $L\subseteq\Fq^\times$, every evaluation $V_t(\xi)$ also lies in $\Fq$, proving the first statement.

For the converse, put $c=\Enc_T(t)\in C_0$.  Condition \eqref{eq:table-basefield-radius} and \cref{lem:basefield-decoding} imply that the representing polynomial $P_c$ lies in $\Fq[X]_{<N}$.  By injectivity of evaluation and the definition of the commitment,
\[
  P_c=V_t.
\]
Hence every coefficient $[X^j]V_t=t(j)$ lies in $\Fq$.
\end{proof}

\subsection{Commit once, interpret many relations}
\label{sec:table-commitment-interpretation}

A central reason for using \cref{def:table-commitment} is that later protocols do not change the commitment format.  The word
\[
  w_t=\Enc_T(t)=\ev_L(V_t)
\]
commits simultaneously to the ordered vector $t$, to the Boolean table whose multilinear extension is $\widetilde t$, and to the degree-$<N$ univariate polynomial $V_t$.  Different proof statements merely apply different algebraic functionals to the same coefficients.

The next sections develop these functionals systematically.  In particular:
\begin{enumerate}[label=(\arabic*),leftmargin=*]
  \item multilinear evaluation $\widetilde t(z)$ is a coefficient of $V_t$ multiplied by a public product-form kernel;
  \item an arbitrary inner product $\ip{a}{b}$ is a coefficient of $V_a$ multiplied by the reversal of $V_b$;
  \item pointwise products are reduced to randomly weighted inner products;
  \item lookup, permutation, sparse linear algebra, and R1CS are then assembled from the same table commitments and these basic relations.
\end{enumerate}
No additional encoding is introduced for these applications.  This common representation is what allows heterogeneous relations to be batched later into a single Reed--Solomon proximity test.

\begin{remark}[Commitment versus proof of a relation]
\label{rem:commitment-not-proof}
The map $\Enc_T$ by itself only binds a table to a Reed--Solomon codeword.  It does not prove any evaluation, inner-product, Hadamard, lookup, or R1CS relation.  Those statements require the coefficient-extraction identities and soundness arguments developed in the following sections.  Keeping this distinction explicit will be important when the interactive protocols are later compiled with hash commitments.
\end{remark}
